\bookchapter{Bab 19: Pemecahan masalah dan penyelamatan diri}

\chapterintro{1. Ponsel menjadi lambat dan panas — pertanyaan yang paling sering ditanyakan}

\booksection{Sejujurnya, ponsel mulai lag setelah digunakan sekitar setahun, itu sangat normal. Alasannya belum tentu ponsel rusak, lebih sering adalah “debu menumpuk” — file sampah, proses latar belakang, ruang penyimpanan perlahan terisi.}

Lakukan diagnosis dulu, jangan terburu-buru memperbaiki:

\booksubsection{Periksa ruang penyimpanan}

\begin{enumerate}
  \item \textbf{: Pergi ke Pengaturan Penyimpanan lihat ruang tersisa. Jika kurang dari 10\%, sistem akan menjadi sangat lambat, karena sedang berjuang mengosongkan tempat. Pengalaman saya: pertahankan setidaknya 15\%-20\% ruang tersisa, ponsel baru bisa cukup lancar.}Periksa peringkat konsumsi daya
  \item \textbf{: Pergi ke pengaturan baterai lihat peringkat konsumsi daya 24 jam terakhir. Jika suatu aplikasi mengambil porsi besar tetapi Anda hampir tidak menggunakannya, berarti ia berjalan di latar belakang. Langsung uninstall atau batasi aktivitas latar belakang.}Pikirkan apa yang baru-baru ini Anda lakukan
  \item \textbf{: Apakah menjadi lambat setelah upgrade sistem? (Setelah update versi besar menjadi lag beberapa hari pertama adalah normal, sistem membangun ulang indeks di latar belakang). Atau menjadi lambat setelah memasang aplikasi baru? (Coba uninstall).}Lihat lingkungan penggunaan
  \item \textbf{: Ponsel terpapar sinar matahari langsung, memakai casing tebal sambil mengisi daya dan bermain game, panas berlebih adalah pasti. Suhu tinggi juga menyebabkan CPU otomatis menurunkan frekuensi, Anda merasa ponsel tiba-tiba menjadi lambat. Dinginkan dulu lalu coba lagi.}Metode pengobatan (dari sederhana ke kompleks)
\end{enumerate}

\textbf{: Mulai ulang ponsel (menyelesaikan 60\% masalah) Bersihkan ruang penyimpanan (hapus aplikasi dan video yang tidak berguna) Jika aplikasi tertentu bermasalah bersihkan cache-nya atau uninstall dan pasang ulang Senjata terakhir: cadangkan data lalu reset pabrik.}2. Aplikasi force close, tidak merespons — aplikasi tidak patuh

\booksection{Suatu aplikasi force close saat dibuka, atau tiba-tiba macet tidak bergerak saat digunakan. Alur penanganan seperti dokter mendiagnosis, langkah demi langkah:}

Tutup dan buka kembali dulu

\begin{enumerate}
  \item \textbf{: Geser aplikasi dari latar belakang, lalu buka kembali. Sebagian besar bug sementara langsung selesai dengan ini.}Update aplikasi
  \item \textbf{: Pergi ke toko aplikasi lihat apakah ada update. Terkadang karena versi aplikasi terlalu lama, tidak kompatibel dengan sistem.}Periksa izin
  \item \textbf{: Beberapa aplikasi membutuhkan izin tertentu untuk bekerja normal, tetapi Anda tidak memberikannya. Misalnya aplikasi kamera tidak diberi izin kamera, dibuka langsung force close. Pergi ke pengaturan temukan aplikasi tersebut, berikan semua izin yang seharusnya.}Bersihkan cache
  \item \textbf{: Pergi ke Pengaturan Manajemen Aplikasi temukan aplikasi tersebut Penyimpanan Hapus Cache. Perhatikan itu "Hapus Cache" bukan "Hapus Data"! Menghapus data akan menghapus status login dan pengaturan dalam aplikasi Anda, setara dengan reset pabrik untuk aplikasi.}Uninstall dan pasang ulang
  \item \textbf{: Senjata terakhir. Sebelum uninstall konfirmasikan dulu apakah data aplikasi ini akan hilang — untuk WeChat cadangkan riwayat chat dulu, untuk game konfirmasikan dulu apakah akun terikat ke nomor ponsel.}3. Tidak bisa mengisi daya atau pengisian lambat — mungkin bukan masalah baterai
\end{enumerate}

\booksection{Tidak bisa mengisi daya adalah masalah yang sangat mudah membuat cemas, tetapi sering kali penyebabnya ada pada aksesori, bukan ponsel rusak.}

Periksa port pengisian dulu

\begin{itemize}
  \item \textbf{: Gunakan tusuk gigi atau alat kecil non-konduktif, perlahan mengeluarkan debu dan serat kapas dari port pengisian. Saya sendiri pernah mengalami — sehelai benang dari saku jatuh tepat menghalangi port pengisian, dicolokkan kira-kira kontak buruk, dikeluarkan langsung baik.}Coba ganti kabel dan adaptor pengisi daya
  \item \textbf{: Tingkat kerusakan kabel jauh lebih tinggi dari adaptor pengisi daya, terutama kabel yang sering digulung. Ganti dengan kabel dan adaptor yang pasti tidak bermasalah untuk tes silang, cepat bisa menentukan aksesori mana yang rusak.}Masalah pengisian nirkabel
  \item \textbf{: Jika Anda menggunakan pengisian nirkabel, periksa apakah casing ponsel terlalu tebal, atau ada pelat logam di belakang casing (bracket magnetik dll.). Juga, letakkan ponsel tepat di tengah pad pengisian — miring mungkin tidak bisa mengisi.}Pengisian lambat
  \item \textbf{: Konfirmasikan apakah daya adaptor pengisi daya cukup (iPhone untuk pengisian cepat butuh setidaknya adaptor PD 20W), apakah kabel mendukung pengisian cepat (tidak semua kabel bisa melewatkan arus di atas 3A), apakah ponsel terlalu panas (terlalu panas akan otomatis menurunkan daya untuk melindungi baterai).}4. Tidak bisa terhubung jaringan, Wi-Fi terus putus
\end{itemize}

\booksection{Di luar menggunakan data bagus-baik saja, pulang tidak bisa terhubung Wi-Fi; atau Wi-Fi terhubung tetapi tidak ada internet. Periksa dengan urutan ini:}

Aktifkan/matikan mode pesawat

\begin{enumerate}
  \item \textbf{: Tarik ke bawah Pusat Kontrol, aktifkan mode pesawat tunggu 5 detik lalu matikan. Ini sama dengan membiarkan ponsel mencari jaringan lagi, sangat efektif.}Wi-Fi "lupa" lalu sambungkan ulang
  \item \textbf{: Pergi ke pengaturan Wi-Fi, klik tanda seru di sebelah nama jaringan tersebut (Android adalah roda gigi), pilih "Lupakan Jaringan Ini" atau "Lupa", lalu masukkan kata sandi lagi untuk terhubung.}Mulai ulang router
  \item \textbf{: Router di rumah banyak orang berbulan-bulan tidak pernah dimatikan, lama-lama menjadi tidak stabil. Cabut daya tunggu 30 detik lalu colokkan lagi, sebagian besar masalah Wi-Fi teratasi.}Periksa apakah masalahnya di ponsel
  \item \textbf{: Apakah perangkat lain terhubung ke Wi-Fi yang sama normal? Jika semua normal hanya ponsel Anda yang bermasalah, mungkin ponsel perlu mulai ulang atau reset pengaturan jaringan (Pengaturan Umum Transfer atau Reset Reset Pengaturan Jaringan, tidak akan menghapus data Anda, tetapi akan menghapus semua kata sandi Wi-Fi).}5. Masalah layar: Sentuh tidak responsif, layar berkedip, meninggalkan jejak
\end{enumerate}

\booksection{Sentuh tidak responsif}

\begin{itemize}
  \item \textbf{: Lepaskan tempered glass dulu lihat apakah masalahnya di glass — tempered glass murah setelah lama memang bisa mempengaruhi sentuh. Ada air di layar atau minyak di tangan juga bisa menyebabkan sentuh drift. Jika bukan, mulai ulang lalu coba lagi.}Ada jejak/sisa bayangan di layar (burn-in)
  \item \textbf{: Menampilkan gambar yang sama lama (misalnya terus menyala men-scroll TikTok), akan meninggalkan "sisa bayangan" di layar OLED. Ini bukan layar rusak, hanya sementara. Matikan layar istirahat setengah jam umumnya bisa pulih.}Layar berkedip
  \item \textbf{: Jika berkedip pada kecerahan rendah, mungkin masalah modulasi PWM, lebih sering pada layar OLED. Cari opsi "DC Dimming" atau "Anti-Flicker" di pengaturan, setelah diaktifkan akan jauh lebih baik. Tidak semua ponsel memiliki fungsi ini.}Sebagian area layar sepenuhnya gelap/terang
  \item \textbf{: Ini cukup serius, kemungkinan besar kerusakan fisik layar, perlu ganti layar. Layar dari reparasi non-resmi, warna dan kecerahan biasanya tidak sebaik asli, juga mempengaruhi ketahanan air.}6. Tidak bisa dinyalakan, terjebak di antarmuka Logo — jangan panik dulu
\end{itemize}

\booksection{Situasi ini paling mudah membuat jantung berdebar, tetapi tenang, coba ini dulu:}

Isi daya setengah jam dulu

\begin{itemize}
  \item \textbf{: Sering kali hanya daya benar-benar habis. Colokkan pengisi daya asli, tunggu setidaknya 30 menit, lalu coba nyalakan. Beberapa ponsel setelah daya habis perlu diisi hingga tingkat tertentu baru akan merespons.}Force restart
  \item \textbf{: iPhone full-screen tekan volume naik volume turun tahan tombol daya sampai muncul Logo. Android tahan tombol daya 10-15 detik sampai bergetar.}Terjebak di antarmuka boot terus restart
  \item \textbf{: Situasi ini disebut "White Apple" (iPhone) atau "Stuck Logo" (Android). Jika force restart tidak berhasil, mungkin perlu masuk recovery mode gunakan komputer untuk memperbaiki. iPhone hubungkan ke komputer buka iTunes/Finder, masuk recovery mode lalu pilih "Update" (bukan "Restore"), akan reinstall sistem tetapi mempertahankan data. Android masuk mode Fastboot atau Recovery, kombinasi tombol setiap merek berbeda, disarankan cari model ponsel Anda + "masuk recovery mode".}Jika semua di atas tidak berhasil
  \item \textbf{: Maka hanya bisa meminta bantuan reparasi profesional. Jika data pernah dicadangkan, tidak perlu terlalu khawatir.}7. Masalah suara: Speaker tidak bunyi, suara panggilan kecil
\end{itemize}

\booksection{Speaker tidak ada suara}

\begin{itemize}
  \item \textbf{: Pertama pastikan bukan mode senyap (tombol senyap di sisi iPhone), mode jangan ganggu, atau volume ditekan ke paling rendah. Terdengar bodoh, tetapi situasi ini mencakup 40\%. Kemudian periksa apakah Bluetooth terhubung ke suatu earphone atau speaker.}Saat panggilan lawan tidak mendengar Anda berbicara
  \item \textbf{: Periksa apakah lubang mikrofon tersumbat — lubang kecil di bagian bawah ponsel sangat mudah tersumbat debu. Gunakan sikat lembut untuk membersihkan dengan lembut. Jangan gunakan jarum atau tusuk gigi menusuk keras, akan merusak jaring anti-debu.}Suara serak/suara pecah
  \item \textbf{: Jika disebabkan oleh speaker kemasukan air — jangan gunakan pengering rambut angin panas! Gunakan tisu menyerap air di permukaan, lalu biarkan kering secara alami. Online mengatakan menggunakan aplikasi pengusir air dengan gelombang suara ada efek tertentu (memutar suara frekuensi rendah untuk menggetarkan air keluar), bisa dicoba.}8. Cari reparasi resmi, hindari jebakan tersebut
\end{itemize}

\booksection{Jika semua metode di atas sudah dicoba tetap tidak berhasil, maka perlu mencari reparasi profesional. Tetapi pasar reparasi sangat dalam, mari saya ceritakan jebakan dan pengalaman saya sendiri:}

Cari resmi yang sebenarnya, jangan cari yang “terlihat seperti resmi”

\begin{itemize}
  \item \textbf{: Pergi ke situs web merek (Apple.com.cn, Huawei.com.cn dll.) cari daftar "Penyedia Layanan Resmi". Toko pinggir jalan yang menggantung logo Apple non-resmi, mungkin tidak menggunakan aksesori asli, setelah diperbaiki ketahanan air Anda hilang seketika.}Hal yang harus dilakukan sebelum mengirim reparasi
  \item \textbf{: Cadangkan semua data (jangan berharap data tidak hilang saat diperbaiki), keluar dari semua akun penting (Apple ID, WeChat, Alipay), matikan "Temukan Ponsel Saya". Beberapa tempat reparasi akan meminta Anda memberikan kata sandi kunci layar, hal ini ada risiko privasi — sebaiknya Anda keluar dari akun dulu sebelum mengirim.}Jebakan ganti layar
  \item \textbf{: Layar non-resmi dibagi "layar press asli" (kaca asli diganti), "layar press belakang" (kaca non-asli), "layar rakitan" (semua non-asli). Harga beda beberapa kali, tetapi efek tampilan dan daya tahan juga beda jauh. Jika ponsel masih dalam garansi, jujur gunakan resmi.}Jebakan ganti baterai
  \item \textbf{: Baterai non-asli mungkin kapasitas palsu, dan setelah digunakan setengah tahun melemah lebih parah dari baterai asli dua tahun. Harga ganti baterai resmi Apple tidak jauh lebih mahal dari non-resmi, disarankan gunakan resmi.}Jebakan pemulihan data
  \item \textbf{: Setelah ponsel kemasukan air/jatuh rusak, pergi reparasi jika pihak tersebut mengatakan "pemulihan data butuh biaya tambahan", hati-hati — beberapa pedagang tidak jujur akan dengan sengaja melebih-lebihkan tingkat kerusakan untuk mengambil lebih banyak uang. Pemulihan data yang sah adalah memberi harga dulu lalu operasi, dan jika pemulihan tidak berhasil tidak dikenakan biaya.}Langkah pertama serbaguna: Mulai ulang ponsel
\end{itemize}

\begin{tcbwarn}{Apa pun masalah yang Anda hadapi — lag, force close, tidak bisa terhubung jaringan, suara abnormal — mulai ulang ponsel dulu. Kalimat ini terdengar seperti menggampangkan, tetapi berdasarkan pengalaman saya sendiri, setidaknya enam dari sepuluh masalah kecil setelah mulai ulang langsung selesai. Ponsel seperti komputer, berjalan lama di memori akan menumpuk berbagai hal berantakan, mulai ulang adalah mengosongkan dan memulai lagi. Jadi, lain kali ponsel bermasalah, jangan terburu mencari "apa yang harus dilakukan", setelah mulai ulang lihat lagi.}
Ingat sekarang: Operasi penyelamat ini harus hafal di luar kepala

\end{tcbwarn}

\begin{tcbpractice}{Force restart}
\begin{enumerate}
  \item \textbf{: iPhone full-screen tekan volume naik volume turun tahan tombol daya sampai muncul Logo. Android tahan tombol daya 10-15 detik sampai bergetar.}Bersihkan port pengisian
  \item \textbf{: Cari alat kecil (tusuk gigi atau penusuk kartu ponsel), perlahan bersihkan debu di port pengisian.}Sambungkan ulang Wi-Fi
  \item \textbf{: Pergi ke pengaturan Wi-Fi "lupa" Wi-Fi Anda saat ini, masukkan kata sandi lagi untuk terhubung.}Aktifkan/matikan mode pesawat
  \item \textbf{: Tarik ke bawah Pusat Kontrol, aktifkan mode pesawat, hitung 5 detik dalam hati, matikan, lihat apakah jaringan pulih.}Screenshot langkah-langkah operasi ini simpan, atau tulis di catatan. Saat menghadapi masalah tidak perlu panik mencari, langsung buka catatan saja.
  \item Bab 20: Mode pengembang dan izin lanjutan
\end{enumerate}

\end{tcbpractice}
